\section{OR Nodes vs.\ Dimensions: What $D_{ff}$ Is Actually Counting}
\label{sec:or-nodes-dimensions}

There is a persistent ambiguity in how to count the components of the implicit
factor graph. The formal construction has one OR gate per factor node, implemented
by one FFN hidden unit. A production transformer has $D_{ff} \approx 14{,}336$
hidden units per layer. Are there 14,336 factor nodes per layer? That seems too
many. This section resolves the ambiguity.

\subsection*{Hidden Units Are Not Factor Nodes}

In the formal construction, one factor node corresponds to one conditional
relationship between two propositions. One FFN hidden unit implements one such
relationship. So in the minimal construction, hidden units and factor nodes are
in one-to-one correspondence.

In a production transformer, this one-to-one correspondence breaks down in two
ways:

\begin{enumerate}
\item \textbf{Multiple units per factor.} A single conditional relationship may
require multiple hidden units to represent accurately with non-sigmoid activations.
ReLU and SwiGLU require more units to approximate the exact log-odds combination
that sigmoid implements in one unit. The factor node is still the unit of
semantics; the hidden units are the implementation.

\item \textbf{Superposed factors.} Because the residual stream represents many
concepts via superposition, a single hidden unit may participate in multiple
factor relationships simultaneously. The unit fires for a direction in residual
stream space that is a superposition of several feature directions. The
contribution it writes is a superposition of several belief updates.
\end{enumerate}

The right level of abstraction is not hidden units but \emph{SAE features}. Each
monosemantic SAE feature corresponds to one factor node in the implicit factor
graph. The number of factor nodes per layer is the number of active SAE features,
not $D_{ff}$.

\subsection*{What $D_{ff}$ Is Counting}

$D_{ff}$ is the \emph{implementation budget} for the OR gates, not the number of
factor nodes. It sets an upper bound on how many distinct conditional relationships
can be represented per layer, but the actual number of semantically distinct
relationships is determined by the SAE feature count, which is typically much
smaller than $D_{ff}$.

The ratio $D_{ff} / D_{\mathrm{model}} \approx 3$--$4$ reflects the implementation
overhead: roughly 3--4 hidden units are needed per residual stream dimension to
handle superposition, non-sigmoid activations, and $k$-ary neighbor chains. The
factor graph has far fewer nodes than $D_{ff}$ suggests; $D_{ff}$ is the engineering
budget for implementing those nodes accurately.

\subsection*{The Dimension Budget}

A cleaner way to think about the implicit factor graph uses the dimension budget
directly. The residual stream has $D_{\mathrm{model}}$ dimensions. Each active
concept occupies a direction in this space. The number of simultaneously active
concepts is bounded by the superposition capacity --- empirically, roughly
$10 \times D_{\mathrm{model}}$ features can be superposed with manageable
interference~\cite{elhage2022superposition}.

For Llama~3 405B with $D_{\mathrm{model}} = 16{,}384$, this gives roughly 160,000
simultaneously active features per token position. Each feature is a factor graph
node. The factor graph has on the order of $W \times 160{,}000 \approx 1.3\mathrm{B}$
nodes for a sequence of length 8192. This is the scale of the implicit Bayesian
network that one forward pass of Llama~3 405B is performing inference over.

\subsection*{The OR/AND Asymmetry Revisited}

With this clarification, the $D_{ff}/H \approx 416$ ratio has a cleaner
interpretation. It is not that there are 416 OR gates for every AND gate. It is
that the OR implementation requires 416 units of hidden dimension for every unit
of attention head dimension. The factor graph itself may be much sparser ---
each node may have only 2--5 neighbors in practice --- but implementing those
sparse connections accurately in the superposed representation requires the
$416\times$ overhead.

The implication for interpretability: finding the factor graph requires working at
the SAE feature level, not the hidden unit level. The factor potentials are encoded
in the relationship between input SAE features (what the $W_1$ rows respond to)
and output SAE features (what the $W_2$ columns write). The $D_{ff}$ hidden units
are the implementation substrate; the SAE features are the semantics.